El argumento comienza con la compatibilidad de prefijos y marcos,
construye el código ternario y determina la bandera mediante \(K\) y el
registro firmado de \(\alpha\). La incidencia marcada selecciona después
el origen reticular y el vecino de Leech. El último paso realiza el álgebra
de vértices y sus automorfismos. En cada etapa se precisará qué estructura
se conserva: combinaciones lineales, soportes, clases discriminantes, forma
bilineal o producto graduado. Las explicaciones de la función de \(K\),
de la normalización y del origen de norma \(54\) acompañan sus respectivas
construcciones en las secciones \ref{exc:bandera} y \ref{exc:leech}.
